\section{Testbed and Measurement Process}
\label{sec:testbed}

\paragraph{Architecture and task.} All experiments use a vanilla
single-head, single-layer DeltaNet with the production
\texttt{chunk\_delta\_rule} kernel \citep{yang2024deltanet}, whose state
update is
$S_t = S_{t-1}\left(I - \beta_t k_t k_t^{\top}\right) + \beta_t v_t
k_t^{\top}$ \citep{schlag2021linear}. The task is synthetic-grammar
associative recall over a pool of 107 trainable entities: within an
episode, each queried entity must be bound to a value through the delta
rule and recovered later through multi-hop composition. Models are trained
for 20{,}000 steps per cell; each (load, seed) pair is one cell.

\paragraph{Readout.} Recovery is graded by \emph{exact continuous} cosine
similarity against the true value vector at threshold $0.9$, written
$\mathrm{rec}@0.9$, never by argmax over a codebook. The distinction is
load-bearing: under argmax decoding a rank-one state can recover on the
order of $d$ associations \citep{nichani2025factual}, so a capacity
measurement that permits codebook lookup measures the codebook, not the
state. The paper's single load-bearing quantity is $h4$, the
$\mathrm{rec}@0.9$ fraction on \emph{held-out} entities and bindings at
hop depth 4, so every reported recovery is a compositional generalization,
not a memorized pair.

\paragraph{The load parameter.} $K$ is the number of independent entity
bindings the state must support simultaneously within one episode;
$K/d_{\mathrm{state}}$ is the capacity ratio. The frontier statistic
$x_0$ is the fitted midpoint of a logistic
$h4(x) = L / (1 + \exp((x - x_0)/w))$ in $x = K/d_{\mathrm{state}}$: the
load at which recall crosses half of its ceiling. $x_0$ is a well-defined
transition location, not the largest load with perfect recall; the fitted
width $w$ reports how sharp the crossing is.

\paragraph{The anchor table and coherence.} A per-entity key-anchoring
table (established in companion work on this testbed as improving held-out
recovery, and treated here as part of the fixed architecture) is the
object whose geometry differs across scales. Its coherence is
$\max\lvert\cos\rvert$ over the table's 107 row vectors. The Welch bound
\citep{welch1974lower} forces this strictly positive whenever
$n = 107 > d_{\mathrm{state}}$, as at $d_{\mathrm{state}}=64$, and permits
exact orthogonality when $n \le d_{\mathrm{state}}$, as at
$d_{\mathrm{state}}=128$ where the constructed table measures
$\max\lvert\cos\rvert = 0.000000$. % evidence: C5
This single structural difference motivates the dose-response test of
Section~\ref{sec:dose}.

\paragraph{Pre-registration and instrument discipline.} Each wave's
hypothesis and outcome semantics were fixed before its cells ran. Three
rules matter downstream. (i) \emph{Degenerate-fit disclosure:} a sigmoid
fitted to flat-at-ceiling data is degenerate by construction; the
registered rule is to report the degenerate bootstrap fraction and refuse
to quote the point estimate as a located transition. (ii)
\emph{Instrument saturation is checked, not assumed:} hop-21 recovery, a
far harder extrapolation than $h4$, is recorded at every wave; a flat
$h4$ result is reportable only when hop-21 is not uniformly $1.0$, ruling
out a globally ceilinged readout. (iii) \emph{Admission gates:} every cell
passes mechanical admissibility checks (convergence of the eval-time
Newton--Schulz block, value-salvage floors, finite loss) before its $h4$
enters any fit; Section~\ref{sec:frontier} reports the one place where the
gate itself, not the measurement, required recalibration.
